%!TEX TS-program = xelatex
%!TEX encoding = UTF-8 Unicode
\documentclass[11pt,a4paper,fontset=fandol]{ctexart}

\usepackage[margin=2.4cm]{geometry}
\usepackage{amsmath}
\usepackage{booktabs}
\usepackage{xcolor}
\usepackage{listings}
\usepackage{hyperref}

\hypersetup{colorlinks=true,linkcolor=blue,urlcolor=blue}

\lstset{
  basicstyle=\ttfamily\small,
  keywordstyle=\color{blue},
  commentstyle=\color{green!45!black},
  numbers=left,
  numberstyle=\tiny,
  frame=single,
  breaklines=true,
  columns=fullflexible
}

\title{题目 2：NEP 势函数 OpenMP 并行优化代码修改说明}
\author{}
\date{2026 年 6 月 6 日}

\begin{document}
\maketitle

\section{修改目标}

本次修改针对 ABACUS 分子动力学模块中的 NEP CPU 势函数计算路径，目标是完成题目
2 所要求的 OpenMP 多线程加速，并保证多线程计算与单线程计算结果一致。

主要修改文件如下：
\begin{itemize}
  \item \texttt{potential/ml/nep/nep\_cpu.cpp}：实现 NEP 力计算路径的 OpenMP 并行；
  \item \texttt{test/nep\_omp\_test.cpp}：增加单线程和多线程结果一致性测试；
  \item \texttt{test/CMakeLists.txt}：注册新增的 NEP OpenMP 单元测试；
  \item \texttt{tools/nep\_omp\_sweep.cpp}：实现不同原子数和线程数的测试程序；
  \item \texttt{tools/run\_problem2\_nep\_sweep.sh}：编译并运行完整测试，保存 CSV 结果。
\end{itemize}

\section{原始代码的问题}

\subsection{可以直接并行的部分}

NEP 代码中的邻居表构建和描述符计算已经以中心原子 \(n_1\) 为循环变量进行并行。
不同线程主要读取共享输入，并写入与 \(n_1\) 对应的独立输出位置，因此不存在明显写冲突。

\subsection{力计算中的数据竞争}

力计算不能简单地在最外层循环前加入
\texttt{\#pragma omp parallel for}。对于一对原子 \(n_1\) 和 \(n_2\)，计算结果会同时更新
两个原子的受力：
\[
F_{n_1} \mathrel{+}= f_{12}, \qquad
F_{n_2} \mathrel{-}= f_{12}.
\]

虽然不同线程负责不同的中心原子 \(n_1\)，但是它们可能同时把同一个原子当作邻居
\(n_2\)。因此多个线程会同时修改相同的 \texttt{force} 和 \texttt{virial} 数组元素，
形成数据竞争，导致结果不稳定或错误。

原始代码中的典型共享写操作如下：

\begin{lstlisting}[language=C++]
g_fx[n1] += f12[0];
g_fx[n2] -= f12[0];
g_virial[n2 + 0 * N] -= r12[0] * f12[0];
\end{lstlisting}

\section{并行设计方案}

\subsection{线程私有缓冲区}

本次修改采用线程私有缓冲区加最终归约的方案。每个 OpenMP 线程分别创建自己的
受力和维里数组：

\begin{lstlisting}[language=C++]
std::vector<double> fx_private(N, 0.0);
std::vector<double> fy_private(N, 0.0);
std::vector<double> fz_private(N, 0.0);
std::vector<double> virial_private(N * 9, 0.0);
\end{lstlisting}

在线程执行邻居循环时，只修改当前线程的私有数组，不修改共享输出数组。完成并行循环
之后，再将各线程的私有结果累加到最终结果。

这种设计具有以下优点：
\begin{itemize}
  \item 避免内层邻居循环中的数据竞争；
  \item 不需要对每一次力更新使用原子操作；
  \item 保持原有力和维里计算公式不变；
  \item 单线程和多线程路径使用相同的核心计算逻辑。
\end{itemize}

\subsection{统一归约辅助函数}

在 \texttt{nep\_cpu.cpp} 的匿名命名空间中新增了
\texttt{reduce\_private\_buffer} 函数：

\begin{lstlisting}[language=C++]
void reduce_private_buffer(
  const std::vector<double>& source,
  double* target)
{
  if (target == nullptr || source.empty()) {
    return;
  }

  for (std::size_t i = 0; i < source.size(); ++i) {
    target[i] += source[i];
  }
}
\end{lstlisting}

该函数负责把一个线程的私有数组累加到共享输出数组中，并处理空指针和空数组情况。

\section{三个力计算函数的具体修改}

\subsection{径向力计算}

修改函数：
\begin{lstlisting}
find_force_radial_small_box
\end{lstlisting}

函数外层加入 OpenMP 并行区域，每个线程建立私有
\texttt{fx}、\texttt{fy}、\texttt{fz} 和 \texttt{virial} 数组。中心原子循环通过
\texttt{omp for schedule(static)} 分配给各线程：

\begin{lstlisting}[language=C++]
#pragma omp parallel
{
  std::vector<double> fx_private(N, 0.0);
  ...

#pragma omp for schedule(static)
  for (int n1 = 0; n1 < N; ++n1) {
    ...
  }
}
\end{lstlisting}

原本对共享数组 \texttt{g\_fx}、\texttt{g\_fy}、\texttt{g\_fz} 和
\texttt{g\_virial} 的写操作，改为写入线程私有指针 \texttt{fx}、\texttt{fy}、
\texttt{fz} 和 \texttt{virial}。

\subsection{角向力计算}

修改函数：
\begin{lstlisting}
find_force_angular_small_box
\end{lstlisting}

角向力计算采用与径向力相同的线程私有缓冲区方案。每个线程还会在栈上保存中心原子
对应的 \texttt{Fp} 和 \texttt{sum\_fxyz} 临时数据，因此这些变量天然为线程私有，
不需要额外同步。

\subsection{ZBL 短程排斥力计算}

修改函数：
\begin{lstlisting}
find_force_ZBL_small_box
\end{lstlisting}

ZBL 路径除了受力和维里，还会更新每个原子的势能 \texttt{g\_pe}。因此该函数额外创建
\texttt{pe\_private}：

\begin{lstlisting}[language=C++]
std::vector<double> pe_private(N, 0.0);
\end{lstlisting}

并行循环结束后，力、维里和势能私有数组统一归约到共享输出。

\subsection{归约过程}

每个并行区域结束前使用临界区完成归约：

\begin{lstlisting}[language=C++]
#pragma omp critical
{
  reduce_private_buffer(fx_private, g_fx);
  reduce_private_buffer(fy_private, g_fy);
  reduce_private_buffer(fz_private, g_fz);
  reduce_private_buffer(virial_private, g_virial);
}
\end{lstlisting}

临界区只用于最终数组累加，不位于计算量较大的邻居循环内部。因此相比在每次原子对更新
时使用 \texttt{omp atomic}，同步次数显著减少。

\section{兼容单线程和非 OpenMP 编译}

所有 OpenMP 代码都由 \texttt{\_OPENMP} 宏保护：

\begin{lstlisting}[language=C++]
#if defined(_OPENMP)
  // OpenMP parallel implementation
#else
  // Original serial pointers
#endif
\end{lstlisting}

当编译器没有启用 OpenMP 时，代码继续直接使用原始输出指针，保持串行行为。因此本次
修改不会破坏非 OpenMP 构建。

\section{正确性测试修改}

\subsection{新增单元测试}

新增文件 \texttt{test/nep\_omp\_test.cpp}。测试使用项目中已有的
\texttt{nep\_hfo2.txt} 模型，构造一个确定性的六原子 Hf/O 体系。

测试步骤如下：
\begin{enumerate}
  \item 调用 \texttt{omp\_set\_num\_threads(1)}，获得单线程参考结果；
  \item 调用 \texttt{omp\_set\_num\_threads(4)}，获得多线程结果；
  \item 分别比较势能、受力和维里数组。
\end{enumerate}

测试容差设置为：
\[
|\Delta E| \le 10^{-12}, \quad
|\Delta F| \le 10^{-10}, \quad
|\Delta V| \le 10^{-10}.
\]

\subsection{CMake 测试接入}

在 \texttt{test/CMakeLists.txt} 中新增测试目标
\texttt{MODULE\_MD\_nep\_omp}，并编译 NEP 所需源文件。模型文件路径通过
\texttt{NEP\_TEST\_MODEL\_FILE} 编译宏传入，使测试不依赖运行时工作目录。

\section{性能测试脚本修改}

\subsection{测试程序}

\texttt{tools/nep\_omp\_sweep.cpp} 实现以下功能：
\begin{itemize}
  \item 根据指定原子数生成规则 Hf/O 体系；
  \item 支持多个原子规模和多个 OpenMP 线程数；
  \item 使用单线程结果作为正确性参考；
  \item 重复调用 \texttt{NEP::compute} 测量平均运行时间；
  \item 输出时间、加速比和最大绝对误差。
\end{itemize}

模型加载在计时循环外完成，因此测量结果主要反映 NEP 计算过程，而不是文件读取时间。

\subsection{运行脚本}

\texttt{tools/run\_problem2\_nep\_sweep.sh} 自动完成：
\begin{enumerate}
  \item 查找支持 OpenMP 的 C++ 编译器；
  \item 使用 \texttt{-O2 -fopenmp} 编译测试程序；
  \item 执行不同规模和线程数组合；
  \item 使用 \texttt{tee} 保存 CSV 结果。
\end{enumerate}

默认运行命令为：

\begin{lstlisting}[language=bash]
cd source/source_md
tools/run_problem2_nep_sweep.sh
\end{lstlisting}

\section{修改结果}

完整测试覆盖了 64、216 和 512 个原子的体系，以及 1、2 和 4 个 OpenMP 线程。结果
显示：
\begin{itemize}
  \item 所有多线程势能结果与单线程结果完全一致；
  \item 受力最大误差约为 \(10^{-15}\)；
  \item 维里最大误差约为 \(10^{-14}\)；
  \item 4 线程运行获得约 \(2.94\) 至 \(3.55\) 倍加速。
\end{itemize}

这些误差属于浮点数累加顺序变化引起的正常舍入误差，说明修改后的代码不存在可观察的
线程安全问题。

\section{开销与可继续优化方向}

当前方案的主要额外开销是线程私有数组，其空间复杂度约为：
\[
O(TN),
\]
其中 \(T\) 为线程数，\(N\) 为原子数。

此外，最终归约使用临界区，线程数很高时可能成为瓶颈。后续可以继续研究：
\begin{itemize}
  \item 将数组分块后进行分层归约；
  \item 复用线程私有缓冲区，减少每次计算时的内存分配；
  \item 使用支持数组归约的 OpenMP 版本；
  \item 结合 MPI 空间分解，减少单进程中的原子规模。
\end{itemize}

\section{总结}

本次修改没有改变 NEP 势函数的物理计算公式，而是重新组织了力、维里和 ZBL 势能的
写入方式。通过线程私有缓冲区消除数据竞争，再通过最终归约合并结果，实现了安全且具有
明显加速效果的 OpenMP 力计算路径。同时新增了单元测试、性能扫描程序和自动运行脚本，
使正确性和性能结果均可重复验证。

\end{document}
